% Responsable: Aldair Herrera (Documentación)
\section{Objetivos}

\subsection{Objetivo general}
Diseñar y construir, a lo largo del periodo SPA 2026--2027, una versión distribuida del
SSDC que permita recibir las hojas de vida de los postulantes, calificarlas con la matriz
de valoración y generar el informe de selección, a partir de los requisitos obtenidos en
las entrevistas al Decano y a la Coordinadora de la carrera.

\subsection{Objetivos específicos}

\begin{table}[H]
\centering
\small
\caption{Objetivos específicos e indicadores}
\begin{tabularx}{\linewidth}{@{}p{0.9cm}XX@{}}
\toprule
\textbf{ID} & \textbf{Objetivo} & \textbf{Indicador de cumplimiento} \\
\midrule
OE-1 & Obtener los requisitos del sistema a partir de entrevistas a quienes realizan hoy la selección. &
Al menos 10 requisitos funcionales y 6 no funcionales, cada uno relacionado con un fragmento de las entrevistas (E1). \\
\midrule
OE-2 & Definir la arquitectura distribuida base del SSDC y registrar sus decisiones. &
Diagramas C4 de nivel 1 y 2 versionados en \texttt{docs/diagramas/} y al menos 2 ADR en \texttt{docs/adr/} (E1). \\
\midrule
OE-3 & Construir un esqueleto de servicios que se comuniquen con el mecanismo elegido. &
Al menos 2 servicios se levantan con un solo comando siguiendo el README desde un clon limpio (E1); todos los servicios del C4 en contenedores (E2). \\
\midrule
OE-4 & Automatizar la calificación con la matriz de valoración. &
En los casos de prueba, el puntaje total que calcula el sistema coincide con el cálculo manual en el 100\,\% de los casos (E4). \\
\midrule
OE-5 & Identificar y mitigar los riesgos técnicos de distribuir el sistema. &
Al menos 5 riesgos con probabilidad, impacto, mitigación y verificación, 3 de ellos derivados de los incidentes de FORM-AC-01 (E1). \\
\bottomrule
\end{tabularx}
\end{table}